\documentclass[11pt,a4paper]{article}

% ──────────────────────────────────────────────────────────────────────────
% AXON-RFC-006 附录 B（实现规约）—— easynet.web_app（中文版）
% Status: 草稿 v0.1（技术设计 + 实现计划）
% Build:  xelatex AXON-RFC-006-B-easynet-webapp.zh-CN.tex
% ──────────────────────────────────────────────────────────────────────────

\usepackage[margin=1in]{geometry}
\usepackage{amsmath,amssymb,amsthm}
\usepackage{xcolor}
\usepackage[colorlinks=true,linkcolor=blue!60!black,citecolor=blue!60!black,urlcolor=blue!60!black]{hyperref}
\usepackage{booktabs}
\usepackage{longtable}
\usepackage{array}
\usepackage{tabularx}
\usepackage{xltabular}
\newcolumntype{L}[1]{>{\raggedright\arraybackslash}p{#1}}
\usepackage{listings}
\usepackage{enumitem}
\usepackage{tcolorbox}
\usepackage{titlesec}
\usepackage{fancyhdr}
\usepackage{parskip}
\usepackage{xspace}

\usepackage{fontspec}
\usepackage{xeCJK}
\setCJKmainfont{PingFang SC}[
    BoldFont={PingFang SC Semibold},
    ItalicFont={PingFang SC Light}
]

\pagestyle{fancy}
\fancyhf{}
\fancyhead[L]{\small AXON-RFC-006 附录 B}
\fancyhead[R]{\small easynet.web\_app —— 技术设计}
\fancyfoot[C]{\thepage}
\renewcommand{\headrulewidth}{0.4pt}

\titleformat{\section}{\Large\bfseries\color{black!85}}{\thesection}{0.7em}{}
\titleformat{\subsection}{\large\bfseries\color{black!75}}{\thesubsection}{0.6em}{}
\titleformat{\subsubsection}{\normalsize\bfseries\color{black!70}}{\thesubsubsection}{0.5em}{}

\definecolor{normativecolor}{RGB}{180,30,30}
\definecolor{derivedcolor}{RGB}{30,90,180}
\definecolor{deferredcolor}{RGB}{120,90,30}
\definecolor{codebg}{RGB}{245,245,248}
\definecolor{codeborder}{RGB}{210,210,220}

\newcommand{\normsv}{\textcolor{normativecolor}{\textbf{[规范性]}}\xspace}
\newcommand{\derived}{\textcolor{derivedcolor}{\textbf{[推论]}}\xspace}
\newcommand{\deferred}{\textcolor{deferredcolor}{\textbf{[暂缓]}}\xspace}

\lstdefinestyle{ealstyle}{
    backgroundcolor=\color{codebg},
    basicstyle=\ttfamily\small,
    frame=single,
    rulecolor=\color{codeborder},
    breaklines=true,
    showstringspaces=false,
    keywordstyle=\color{blue!70!black}\bfseries,
    commentstyle=\color{green!40!black}\itshape,
    stringstyle=\color{red!60!black},
    morekeywords={state_type,state_key,canonical,operational,transition,receipt,view,principal,agent,owner_signature,attempt_id,transition_id,routes_manifest,api_abilities,fn,let,struct,impl,pub,async,await,location,server,proxy_pass},
    columns=fullflexible,
    keepspaces=true,
    aboveskip=0.6em,
    belowskip=0.6em,
}
\lstset{style=ealstyle}

\newtheorem{theorem}{定理}
\newtheorem{lemma}[theorem]{引理}
\newtheorem{corollary}[theorem]{推论}
\theoremstyle{definition}
\newtheorem{definition}{定义}

\newtcolorbox{invariant}[1][不变量]{
    colback=red!4!white,
    colframe=red!50!black,
    fonttitle=\bfseries,
    title={不变量 --- #1},
    sharp corners=south,
    boxrule=0.6pt,
}
\newtcolorbox{notebox}[1][]{
    colback=blue!4!white,
    colframe=blue!50!black,
    fonttitle=\bfseries,
    title={说明},
    sharp corners=south,
    boxrule=0.6pt,
    #1
}
\newtcolorbox{warningbox}[1][]{
    colback=orange!6!white,
    colframe=orange!70!black,
    fonttitle=\bfseries,
    title={过渡期偏差},
    sharp corners=south,
    boxrule=0.6pt,
    #1
}
\newtcolorbox{decisionbox}[1][]{
    colback=purple!4!white,
    colframe=purple!60!black,
    fonttitle=\bfseries,
    title={待决策},
    sharp corners=south,
    boxrule=0.6pt,
    #1
}

\title{
    \vspace{-2em}
    \textbf{\Large AXON-RFC-006 附录 B：easynet.web\_app}\\[0.3em]
    \large 由 agent 部署的全栈应用：技术设计、构建管线、Hub 路由
    与实现计划
}
\author{
    胡思蓝 \\ \small 新加坡国立大学
}
\date{草稿 v0.1，\today}

\begin{document}
\maketitle

\begin{abstract}
\noindent
本附录规定 \texttt{easynet.web\_app}：用于支撑完整现代 web 部署的
state\_type。一个 agent（典型情形为 Claude Code）在其 workspace 内
撰写真实的 React/Vue/Next 风格前端项目，通过构建产生静态 \texttt{dist/}，
该 \texttt{dist/} 与一个由 EasyNet ability 直接充当的后端一起被服务。
Hub 扮演\emph{有状态版的 nginx} 角色：服务前端 bundle、把
\texttt{POST /api/<verb>} 翻译成 owner agent 上的 ability 调用、把
WebSocket 订阅代理到 ability stream，并通过
\texttt{/\_\_easynet/agent} 暴露同一应用对机器可执行的视图。本附录
覆盖现代项目模型、canonical / runtime 状态拆分、长任务构建管线、
API ability 映射、Hub 路由矩阵（含适用于 SPA 的 CSP 与缓存规则）、
让整个应用对邻居 agent 可发现的 agent\_view 形态、在 \texttt{EasyNet-Cli}
中落地所需的五个 PR 切片，以及一个端到端的 worked example——agent 用
Vite + React 写一个购物前端、声明 API ability、构建、发布、上线，
并被审计的全过程。
\end{abstract}

\tableofcontents
\vspace{1em}
\hrule
\vspace{1em}

% ──────────────────────────────────────────────────────────────────────────
\section{相对 pages 与 web 网络历史的位置}
\label{sec:position}

\derived{} web 从一个静态文档媒介演变成完整应用平台。EasyNet 有状态
web 层走相同的弧，分两份附录：

\begin{center}
\renewcommand{\arraystretch}{1.25}
\begin{tabularx}{\linewidth}{l X}
\toprule
\textbf{附录} & \textbf{建模对象} \\
\midrule
A（\texttt{easynet.page}）   & 单一文档。作者写 HTML / Markdown；Hub 服务它。最简参考——单一规范状态、无 runtime 叠加、无构建步骤、无 API 面。 \\
B（\texttt{easynet.web\_app}）& 完整应用。作者写真实前端项目，构建管线产生静态 bundle，应用的 API endpoint \emph{即} EasyNet ability。Hub 在统一 URL 空间内服务前端 + API + WebSocket + agent\_view。 \\
\bottomrule
\end{tabularx}
\end{center}

\noindent 两个都保留。Pages 仍是协议规范核最简化的验证；web\_app
是第一个真正调用完整压力测试面（markdown 配套 \S B.1--B.8）的状态类型。
只交付 pages 而不交付 web\_app 的实现，等于解决了有状态 EasyNet 简单的
那一半；本附录规定的是难的那一半。

\textbf{为什么这不是 nginx。} nginx 只服务一个静态目录、并把请求反向
代理到独立部署的后端。后端对 nginx 是黑盒——它不理解后端状态，不知道
后端有哪些 endpoint，也无法对一个通过结构化发现机制询问"调用方在这里
能做什么"的邻居作答。本附录规定的 Hub 服务的字节与 nginx 一样，但其
路由决策派生自 state object 的规范状态、其 API endpoint 是和其他每个
EasyNet 调用方共享同一注册表的、可内省的 ability，并且双视图契约让
同一 URL 空间同时承载人类内容和对邻居 agent 可执行的整个应用地图。
nginx 是 HTTP 管道；Hub 是恰好在外边界讲 HTTP 的状态投影层。

% ──────────────────────────────────────────────────────────────────────────
\section{目标与非目标}
\label{sec:goals}

\subsection{目标}

\normsv{} 一个符合本附录的实现：

\begin{enumerate}[leftmargin=2.2em,itemsep=0.2em]
    \item 让 agent 在自己的 workspace 中用任意工具链（Vite、Next.js、
    Remix 等）撰写真实前端项目；EasyNet 不发明项目结构。
    \item 提供长任务规范 transition \texttt{webapp.build}，运行项目
    声明的构建命令、把产出的 \texttt{dist/}（或等价物）以一份不可变的
    内容寻址 bundle 形式捕获，并按成功/失败推进 canonical state。
    \item 提供 runtime 叠加 transition（\texttt{webapp.serve}、
    \texttt{webapp.stop}）绑定本地端口并开始服务——关键一点：
    不改 canonical state。
    \item 在 Hub 上把 HTTP 请求按三条正交通道翻译为 EasyNet 操作：
    静态 bundle 拉取、API ability 调用、WebSocket ability 订阅。
    \item 暴露单一的 \texttt{agent\_view} URL
    （\texttt{/\_\_easynet/agent}）返回应用完整可发现性面：路由表、
    API endpoint、对 transition 的 \texttt{ability\_surface}、构建
    元数据、回执历史指针、EAL 提示。
    \item 上述全部能在五个 PR 切片（\S\ref{sec:impl-plan}）中落地，
    每个切片可独立审阅。
\end{enumerate}

\subsection{非目标}

\derived{} 不在本附录范围：

\begin{itemize}[leftmargin=2.2em,itemsep=0.18em]
    \item \textbf{TLS 终结、HTTP/2 协商、多实例 Hub 集群。}
    属于运维，不属于协议。
    \item \textbf{选定的前端框架。} 附录以 React/Vite 作 worked example
    是因为 Claude Code 默认使用它；协议本身与框架无关。
    \item \textbf{需要每请求保持进程的服务端渲染。} Phase 1 支持
    静态 + Hub 路由 API + WebSocket。需要每个请求都跑 Node.js 进程的
    SSR 推迟到 RFC-006-B v2。
    \item \textbf{跨 realm 或跨节点的应用联邦。} Phase 1：一个 web\_app
    属于唯一 agent、运行在唯一节点。
    \item \textbf{已认证人类（\texttt{principal/human-auth}）。}
    Phase 1 仅支持 \texttt{principal/human-anon}；应用自己设的
    cookie 透传；EasyNet 不签发 cookie。
\end{itemize}

% ──────────────────────────────────────────────────────────────────────────
\section{现代项目模型}
\label{sec:project}

\subsection{workspace 中存放什么}

\normsv{} state object 在磁盘上的真值来源是\emph{一个普通的项目目录}，
归 agent 所有。具体：

\begin{lstlisting}
<workspace>/web_apps/<slug>/
├── easynet.app.toml          # EasyNet 侧的声明
├── package.json              # agent 工具链所有
├── package-lock.json         # 同上
├── src/                      # 框架风格源代码
├── public/
├── vite.config.ts            # 或 next.config.js 等
└── （可选）框架需要的任何其他文件
\end{lstlisting}

EasyNet \textbf{只拥有} \texttt{easynet.app.toml}；其余属于框架。当 agent
运行 \texttt{webapp.build} 时，\texttt{easynet.app.toml} 中声明的命令在
该目录中执行；它在所声明的 \texttt{dist\_dir} 中产出的任何东西都会作为
构建产物被捕获。

\subsection{easynet.app.toml}

\normsv{} EasyNet 用以理解 web\_app 的唯一文件：

\begin{lstlisting}
# easynet.app.toml —— EasyNet 编写并阅读的唯一文件。
# 项目其余部分属于框架。

[manifest]
name         = "shop"
description  = "由 agent 构建的小型陶瓷店。"
visibility   = "public"            # private | realm | public

[build]
command      = "npm install && npm run build"
dist_dir     = "dist"              # 相对项目根
build_timeout_secs = 600
node_version = "20"                # 建议性；执行器可通过稳定镜像注册表
                                   # 钉版本

[serve]
# 调用 webapp.serve 时，运行时绑定一个端口并启动一个进程，其唯一职责
# 是为 Hub 反向代理暴露 API + WS；静态 bundle 从 blob 存储被服务，
# 不从这个进程被服务。
mode         = "static-only"
                # "static-only" : 无后端进程；API 完全由 Hub 发起的
                #                 ability 调用满足。Phase 1 的默认与
                #                 生产形态。
                # "process"      : 执行器启动后端进程；其端口出现在
                #                 runtime 叠加状态中。预留给 v2。

[api]
# 前端\emph{允许}通过 Hub 的 /api/<verb> 路由调用的 ability。每条是
# owner agent 已注册的全限定 ability 名。Hub 会拒绝把 verb 不在此列表
# 的 /api/<verb> 请求翻译；这是 agent 已经显式授权的攻击面。
abilities = [
  "shop.list_items",
  "shop.checkout",
  "shop.order_status"
]

[routes]
# nginx 风格的路由规则。顺序重要。每条规则有 match 模式与 action。
rules = [
  { match = "exact:/",                      action = "static:dist/index.html" },
  { match = "prefix:/assets/",              action = "static:dist/assets/" },
  { match = "regex:^/api/([a-z][a-z0-9_]*)$",
    action = "api:$1",
    methods = ["GET","POST"] },
  { match = "prefix:/ws/",                  action = "ws:$path" },
  { match = "exact:/__easynet/agent",       action = "agent_view" },
  { match = "exact:/__easynet/health",      action = "health" },
  { match = "exact:/__easynet/receipts",    action = "receipts" },
  { match = "fallback",                     action = "spa:dist/index.html" }
]

[csp]
# 收紧或放宽默认值。默认值适用于普通 SPA；需要 'unsafe-inline' 的
# 框架必须显式选择。
extra_script_src = []   # 例如内联引导脚本的 ["'sha256-...'"]
extra_style_src  = ["'unsafe-inline'"]
extra_connect_src= ["self"]
\end{lstlisting}

文件小、声明式，且是参与 canonical state 哈希的唯一对象（与源树和构建
产物哈希并列）。最小应用三个段（\texttt{[manifest]}、\texttt{[build]}、
\texttt{[api]}）即可；routes 默认走上面的 SPA 标准布局。

\subsection{为什么 API 面必须声明而不是自动检测}

\normsv{} 一个 web\_app 的 API 面是\emph{声明}的，不是\emph{发现}的。
当 \texttt{<verb>} 不在 \texttt{[api].abilities} 里时，Hub 拒绝把
\texttt{/api/<verb>} 翻译成调用。两条理由：

\begin{itemize}[leftmargin=2.2em,itemsep=0.18em]
    \item \textbf{审计清晰。} 回执日志精准记录公开 web\_app 暴露的
    ability。审计者在 canonical state 里看到该面，而不是在运行时内省。
    \item \textbf{攻击面收敛。} owner agent 可能托管很多 ability；
    只有显式列出的能从公开 Hub 触达。agent 注册的内部用 ability
    在 web 上不可见。
\end{itemize}

\begin{invariant}[WB-INV-1]
当 \texttt{<verb>} 不在 web\_app 当前 Built canonical state 的
\texttt{api\_abilities} 列表中时，Hub \textbf{必须}拒绝把
\texttt{/api/<verb>} 翻译为 ability 调用。拒绝必须返 HTTP 404
（不是 403）以避免确认内部 ability 的存在。这是身份层 TR-INV-12
建立的信任边界纪律在应用层的实例化。
\end{invariant}

% ──────────────────────────────────────────────────────────────────────────
\section{状态机}
\label{sec:state-machine}

\subsection{state\_type 与 state\_key}

\normsv{}

\begin{lstlisting}
state_type:  easynet.web_app

state_key:
    (realm: string, owner_agent_uri: string, app_slug: string)

  realm           逻辑命名空间；与 pages 同惯例。
  owner_agent_uri 规范权限。即为本 app 每个 CanonicalReceipt 签名的
                  agent。按主文 §1.3 嵌在 key 内（不是字段）。
  app_slug        owner 命名空间内的 kebab-case 标识符
                  (^[a-z0-9][a-z0-9-]{0,62}[a-z0-9]$)。

URA 形态：
  easynet:///r/<realm>/agent/<owner>/web_app/<slug>
\end{lstlisting}

\subsection{规范字段}

\normsv{} 按 CSH-INV-2（主文 \S2.5）大体积值都按内容哈希持有：

\begin{lstlisting}
canonical_fields(web_app) = {
    manifest_hash:       bytes(32),  // 规范化后 [manifest] 段的
                                     //   SHA-256 (name, description,
                                     //   visibility)
    build_config_hash:   bytes(32),  // 规范化后 [build] 段
    api_surface_hash:    bytes(32),  // 规范化后 [api].abilities
                                     //   （已排序）
    routes_manifest_hash:bytes(32),  // 规范化后 [routes].rules
                                     //   （按声明顺序）
    csp_hash:            bytes(32),  // 规范化后 [csp] 段
    source_snapshot_hash:bytes(32),  // 最近一次构建源树的哈希
                                     //   首次构建前为 null
    build_artifact_hash: bytes(32),  // 上次成功构建产生的 dist 树
                                     //   哈希。首次 Built 前为 null。
                                     //   delete 时清零。
    state_value:         string,     // §4.3
    visibility:          string,     // 镜像 manifest.visibility
                                     //   便于快速 scope check
    version:             uint64
}
\end{lstlisting}

\textbf{canonicality 规则。} 上述每个字段在 schema 中均标
\texttt{canonical:true}；schema 驱动的投影 $\pi$ 恰按上述 key 产出投影。

\subsection{规范状态空间}

\begin{center}
\renewcommand{\arraystretch}{1.25}
\begin{tabularx}{\linewidth}{l X}
\toprule
\textbf{state\_value} & \textbf{含义} \\
\midrule
\texttt{Created}   & 已注册 easynet.app.toml；尚未尝试构建。 \\
\texttt{Building}  & 一次 \texttt{webapp.build} 在进行中。瞬态。
                     仅通过 OperationalReceipt 可见；从不作为终态
                     canonical state。 \\
\texttt{Built}     & 已存在至少一次成功构建。
                     \texttt{build\_artifact\_hash} 非空。可被 Served。 \\
\texttt{Failed}    & 最近一次构建失败。owner 可迭代后重试。 \\
\texttt{Deleted}   & 终态；\texttt{build\_artifact\_hash} 清零，blob
                     引用计数递减。 \\
\bottomrule
\end{tabularx}
\end{center}

\subsection{Runtime 叠加状态空间}

\begin{center}
\renewcommand{\arraystretch}{1.25}
\begin{tabularx}{\linewidth}{l X}
\toprule
\textbf{runtime} & \textbf{含义} \\
\midrule
\texttt{Stopped}  & 此 app 没有附着的服务叠加。 \\
\texttt{Starting} & \texttt{webapp.serve} 已被调用；Hub 正在把 slug
                    绑定到其路由。 \\
\texttt{Running}  & Hub 正主动为此 app 翻译 \texttt{/<slug>/...} 请求。 \\
\texttt{Crashed}  & Hub 丢失了映射（进程重启、配置重载错误等）；
                    canonical state 未被触动。可通过 \texttt{webapp.serve}
                    恢复。 \\
\bottomrule
\end{tabularx}
\end{center}

canonical $\times$ runtime 交叉积矩阵在结构上与 markdown 配套
\S B.4.3 完全一致：仅 \texttt{Built} 允许非 Stopped 的 runtime 叠加。

\subsection{Transition 列表}

\begin{center}
\renewcommand{\arraystretch}{1.2}
\begin{tabularx}{\linewidth}{l l X}
\toprule
\textbf{Transition} & \textbf{类} & \textbf{说明} \\
\midrule
\texttt{webapp.create@v1}     & canonical   & $\emptyset \to$ Created。铸造 state\_key，落 manifest/build\_config/routes/api/csp 哈希。 \\
\texttt{webapp.update\_config@v1}
                              & canonical   & Created/Built/Failed $\to$ 同。替换声明配置（manifest, build, api, routes, csp）。\emph{不}触发构建。 \\
\texttt{webapp.build@v1}      & canonical   & Created/Built/Failed $\to$ Building（瞬态）$\to$ Built|Failed。长任务。见 \S\ref{sec:build-pipeline}。 \\
\texttt{webapp.delete@v1}     & canonical   & * $\to$ Deleted。隐式触发 \texttt{webapp.stop} 副作用。 \\
\texttt{webapp.serve@v1}      & operational & Stopped/Crashed $\to$ Starting $\to$ Running。要求 canonical = Built。 \\
\texttt{webapp.stop@v1}       & operational & Starting/Running/Crashed $\to$ Stopped。 \\
\texttt{webapp.get@v1}        & query       & 返回 canonical + runtime 快照。 \\
\texttt{webapp.list@v1}       & query       & 列出某 agent 拥有的 web\_app。 \\
\bottomrule
\end{tabularx}
\end{center}

% ──────────────────────────────────────────────────────────────────────────
\section{构建管线}
\label{sec:build-pipeline}

\subsection{webapp.build@v1 步骤}

\normsv{} \texttt{webapp.build} 时执行器：

\begin{enumerate}[leftmargin=2.2em,itemsep=0.2em]
    \item 读取 \texttt{easynet.app.toml} 并校验。
    \item 计算 \texttt{source\_snapshot\_hash}：项目目录的 Merkle 哈希
    （排除 \texttt{node\_modules/}、\texttt{dist/}、\texttt{.git/}
    以及框架 \texttt{.gitignore} 排除的任何路径）。哈希取自一个排序后
    的 \texttt{(relative\_path, sha256(file\_bytes))} 对列表，已 JCS
    规范化。
    \item 发出 \texttt{build.started} \emph{OperationalReceipt}，含
    源快照哈希、声明命令、预期超时。\emph{非规范}。
    \item 在项目目录内 spawn 构建命令。stdout / stderr 被三通到流式
    日志；按周期发出 \texttt{build.progress} OperationalReceipt，含
    日志累积哈希。可丢失。
    \item 退出时：
    \begin{itemize}[leftmargin=1.5em]
        \item 状态码 0 且 \texttt{dist\_dir} 非空：遍历
        \texttt{dist\_dir}；按 \texttt{(relative\_path, sha256(file\_bytes))}
        对计算 \texttt{build\_artifact\_hash} Merkle 哈希；把整棵
        dist 树作为单个 \texttt{TreeBlob}（一份按 JCS 规范化的
        manifest，列出文件路径与各自 blob 哈希；文件 blob 作为普通
        内容寻址 blob 单独存储）写到 blob 存储。
        \item 状态码非零或 \texttt{dist\_dir} 空：构建失败原因
        payload（最后 64KB 日志 + 退出状态码）并存其哈希。
    \end{itemize}
    \item 发出唯一一条 \emph{CanonicalReceipt}：成功为
    \texttt{build.completed}，失败为 \texttt{build.failed}。这是该
    transition 的唯一规范推进。\texttt{post\_state\_hash} 反映 Built
    或 Failed；\texttt{state\_version} +1。
\end{enumerate}

\textbf{canonicality 纪律。} 瞬态 \texttt{Building}、流式进度、构建
工具 stdout 都是观测。它们绝不进入规范哈希。仅有规范回执的验证者
看到应用一步从 Created $\to$ Built（或 Failed）——这正确：规范事实是
产物，不是过程。

\subsection{TreeBlob 格式}

\normsv{} TreeBlob 是目录树的内容寻址表示。它包含：

\begin{lstlisting}
TreeBlob (JCS 规范化 JSON)
{
  "version": 1,
  "entries": [
    { "path": "index.html",
      "size": 2148,
      "blob": "<file 字节的 sha256>",
      "mode": "0644" },
    { "path": "assets/index-DfX9k.js",
      "size": 187321,
      "blob": "<sha256>",
      "mode": "0644" },
    ...
  ]
}
\end{lstlisting}

树的哈希是 \texttt{sha256(jcs(TreeBlob))}；这就是
\texttt{build\_artifact\_hash}。每个独立文件也是内容寻址 blob。Hub
通过先在 TreeBlob 中按 \texttt{path lookup}、再
\texttt{blob.get(<file-blob-hash>)} 拉取文件；多一次间接但让 blob
存储在多次构建间、跨多个应用间去重相同字节。

\textbf{为什么不直接用 git tree hash。} git 的 tree hash 也是内容寻址，
但用 git 特有方式编码 mode 位与 entry 类型。我们选 JCS 扁平格式以让
哈希函数与 EasyNet 其他地方完全一致（SHA-256 over JCS 字节），避免
出现第二种规范化机制。

\subsection{构建执行器}

\normsv{} Phase 1 交付一个\emph{本地}构建执行器：构建在拥有 agent 的
EasyNet 守护进程的子进程中运行。owner 直接为产生的 CanonicalReceipt
签名。委托构建（源在 agent A，构建在 agent B 的硬件运行）属于未来：
协议形态已就绪（主文 \S5.3），但执行器尚未实现。

\begin{decisionbox}
\textbf{待决策：构建沙箱。} 构建子进程具备 EasyNet 守护进程的进程权限。
恶意构建脚本可能读取守护进程托管的其他 agent 状态。Phase 1 缓解：
构建在项目目录内运行，限制 \texttt{PATH}，仅继承 \texttt{NODE\_VERSION}
与 \texttt{HOME} 等少量环境变量，且更严的 \texttt{ulimit}。完整沙箱
（seccomp / pledge / Docker）属于 RFC-006-B v2。
\end{decisionbox}

% ──────────────────────────────────────────────────────────────────────────
\section{Hub 路由模型}
\label{sec:hub-routing}

Hub 绑定单端口（Phase 1 默认 \texttt{127.0.0.1:8787}）。每个进入的
HTTP 请求按七步处理：

\begin{enumerate}[leftmargin=2.2em,itemsep=0.18em]
    \item \textbf{caller 铸造。} 按主文 \S5.4 构 envelope：
    \texttt{caller = principal/human-anon/<ip-hash>/<sid>}；
    \texttt{caller\_context} 按主文 \S3.6 / TR-INV-13。
    \item \textbf{path 解析。} path 是其中一种：
    \begin{itemize}[leftmargin=1.5em]
        \item \texttt{/<realm>/<owner>/<slug>/...} —— 规范形态。
        \item 配置好的别名域名下任意路径（如 \texttt{shop.alice.example}
        映射到一组特定 (\texttt{realm, owner, slug})）。
    \end{itemize}
    \item \textbf{state 解析。} 在 state 存储中查 web\_app。若不在 Built，
    按 \S\ref{sec:hub-by-state} 返。若 \texttt{visibility} 拒绝 caller，
    返 404。
    \item \textbf{route 匹配。} 按声明顺序在 web\_app 的
    \texttt{[routes].rules} 中跑 path 尾段直到一条命中。
    \item \textbf{action 分派。} 按命中规则的 action 分派
    （static / api / ws / agent\_view / health / receipts / spa fallback）。
    见 \S\ref{sec:hub-actions}。
    \item \textbf{响应构建。} 按 action 设置 CSP、ETag、缓存头、
    content-type。
    \item \textbf{OperationalReceipt。} 发出一条 OperationalReceipt，
    含 caller、caller\_context、命中规则、action、响应状态。
    \emph{非规范}。
\end{enumerate}

\subsection{按状态划分的 Hub 行为矩阵}
\label{sec:hub-by-state}

\begin{center}
\renewcommand{\arraystretch}{1.25}
\begin{tabularx}{\linewidth}{l l X}
\toprule
\textbf{Canonical} & \textbf{Runtime} & \textbf{Hub 响应} \\
\midrule
\texttt{Created}   & *           & 503 "未构建"。\texttt{Retry-After: 60}。 \\
\texttt{Building}  & *           & 503 + 通过 SSE 流式发送 build.progress 回执的小进度页。 \\
\texttt{Built}     & \texttt{Stopped}    & 503 "未在服务"。 \\
\texttt{Built}     & \texttt{Starting}   & 503 "启动中"，\texttt{Retry-After: 2}。 \\
\texttt{Built}     & \texttt{Running}    & 完整路由矩阵如下。 \\
\texttt{Built}     & \texttt{Crashed}    & 503 "可恢复失败" + reason hash；\texttt{Retry-After: 5}。 \\
\texttt{Failed}    & *           & 503 "上次构建失败" + reason hash。 \\
\texttt{Deleted}   & *           & 404。 \\
\bottomrule
\end{tabularx}
\end{center}

\subsection{Action 分派}
\label{sec:hub-actions}

\paragraph{static:<path>}
在 web\_app 的 \texttt{build\_artifact\_hash} TreeBlob 内解析
\texttt{<path>}。拉取 file blob。响应：
\begin{lstlisting}
HTTP/1.1 200 OK
Content-Type: <由扩展名推断；走允许列表映射>
Content-Length: <size>
ETag: "<base64url(blob_hash[..16])>"
Cache-Control: public, max-age=31536000, immutable
                                # 仅对 /assets/*（带哈希文件名）；
                                # 否则：max-age=60, must-revalidate
Content-Security-Policy: <由 app 的 csp_hash 构建>
Referrer-Policy: no-referrer
X-Content-Type-Options: nosniff
\end{lstlisting}

\texttt{immutable} 缓存策略对带哈希文件名的资源安全（Vite/Webpack 输出
的资源文件名内嵌内容哈希）；\texttt{index.html} 等无哈希文件用短的
\texttt{must-revalidate} 窗口。

\paragraph{api:<verb>}
\texttt{<verb>} 是 route 规则的 regex 捕获。Hub：
\begin{enumerate}[leftmargin=2em,itemsep=0.15em]
    \item 校验 \texttt{<owner>.<verb>} 在 \texttt{api\_surface\_hash}
    对应展开的能力列表中。拒绝：HTTP 404（WB-INV-1）。
    \item 把请求 body（Phase 1 大小上限 1 MiB）作 JSON 读取。该 JSON
    \emph{即} ability 的 args，无 envelope 包裹。\texttt{GET} 请求用
    query string；Hub 把 query string 解析为对象。
    \item 用此 caller、caller\_context、解析后的 args 构 EasyNet
    invocation。
    \item 调用 \texttt{<owner>.<verb>}。
    \item 把结果映射到 HTTP：
    \begin{itemize}[leftmargin=1.5em]
        \item \texttt{Succeeded}：200，body = 回执 result body，
        content-type \texttt{application/json}。
        \item 带结构化错误的 \texttt{Failed} 回执：校验类 400，否则 500；
        body = \texttt{\{ error\_code, error\_message, receipt\_id \}}。
        \item 超时：504；含运维超时回执的 receipt\_id。
    \end{itemize}
\end{enumerate}

\textbf{幂等性。} \texttt{GET} 与 \texttt{POST} 都允许；幂等性是 ability
自己的责任，不归 Hub。Hub 不强制 GET 不产 CanonicalReceipt——它是 EasyNet
ability，按其语义决定是否要发。前端选动词；ability 作者在 schema 中
记录幂等约定。

\paragraph{ws:<path>}
Hub 把连接升级为 WebSocket 并绑到一个 ability \emph{stream}（已有的
EasyNet stream 订阅机制）。客户端首条 WebSocket 消息必须是 JSON
\texttt{\{ "ability": "<verb>", "args": \{...\} \}}；之后 Hub 打开
\texttt{<owner>.<verb>} 流并把每个 emit 事件作为 WebSocket 文本帧
转发。首条之后客户端消息回填给 stream 的输入通道。WebSocket 关闭
映射为 stream 关闭。

\paragraph{agent\_view}
返回 \S\ref{sec:agent-view} 规定的 JSON。

\paragraph{health}
返回 \texttt{200 \{"state":"...","runtime":"..."\}}。

\paragraph{receipts}
返回该 state\_key 的规范回执历史，可由 \texttt{?from=<version>\&to=<version>}
约束。

\paragraph{spa:<fallback-path>}
作为最后的 \texttt{fallback} 规则。等价于
\texttt{static:<fallback-path>} 但带 \texttt{Cache-Control: no-cache}，
让 SPA 外壳始终新鲜，而其加载的 hashed 资源不可变。

% ──────────────────────────────────────────────────────────────────────────
\section{agent\_view：整应用可发现性}
\label{sec:agent-view}

\normsv{} 一个 Built/Running web\_app 的 \texttt{agent\_view}：

\begin{lstlisting}
GET /<realm>/<owner>/<slug>/__easynet/agent

200 OK
Content-Type: application/json

{
  "state_type":  "easynet.web_app",
  "state_key":   { "realm":"default","owner":"alice","slug":"shop" },
  "ura":         "easynet:///r/default/agent/alice/web_app/shop",

  "canonical": {
    "state":              "Built",
    "manifest":           { "name":"shop",
                            "description":"由 agent 构建的小型陶瓷店。",
                            "visibility":"public" },
    "build_artifact_hash":"0x4d28ee17...",
    "version":            7,
    "etag":               "TSjuFw0pAvjpAQ=="
  },

  "runtime": {
    "state":     "Running",
    "started_at":"2026-04-28T09:11:02Z",
    "host":      "https://shop.alice.example"
                  // 未配置别名域名时为 null
  },

  "ability_surface": [   // 当前可执行的 transition
    "webapp.update_config@v1",
    "webapp.build@v1",
    "webapp.stop@v1",
    "webapp.delete@v1"
  ],

  "api": {               // 前端可经 /api/* 调用的内容
    "endpoints": [
      { "verb":"shop.list_items",
        "http": "GET  /api/list_items",
        "schema_url": "/__easynet/schema/shop.list_items" },
      { "verb":"shop.checkout",
        "http": "POST /api/checkout",
        "schema_url": "/__easynet/schema/shop.checkout" },
      { "verb":"shop.order_status",
        "http": "GET  /api/order_status",
        "schema_url": "/__easynet/schema/shop.order_status" }
    ]
  },

  "routes": [
    { "match":"exact:/",         "action":"static:dist/index.html" },
    { "match":"prefix:/assets/", "action":"static:dist/assets/" },
    { "match":"regex:^/api/...","action":"api:$1" },
    { "match":"prefix:/ws/",    "action":"ws:$path" },
    { "match":"fallback",       "action":"spa:dist/index.html" }
  ],

  "build_meta": {
    "command":          "npm install && npm run build",
    "dist_dir":         "dist",
    "node_version":     "20",
    "last_build_at":    "2026-04-28T09:09:51Z",
    "build_duration_s": 38.2,
    "framework_hint":   "vite-react"   // 尽力嗅探
  },

  "history": {
    "versions":                    7,
    "latest_canonical_receipt_id": "r_c3d4...",
    "receipt_log_ura":
       "easynet:///r/default/agent/alice/web_app/shop/receipts"
  },

  "eal_hint":
    "call shop.checkout on alice with { item_id: \"mug-12oz\", qty: 1 } timeout 30"
}
\end{lstlisting}

\textbf{为什么这是 killer feature。} 邻居 agent 读到此 URL，一份文档里
就拿到了它需要的一切：

\begin{itemize}[leftmargin=2.2em,itemsep=0.15em]
    \item 知道这是什么应用、它处于何种状态；
    \item 精确知道前端可以调用哪些 API verb（含路径、方法、schema 指针）；
    \item 知道哪些 transition 推进该应用（\texttt{ability\_surface}）；
    \item 既能直接经 EasyNet 调用任何一个 verb，也能经公开
    \texttt{/api/<verb>} 调用——两条路径触达同一个 ability；
    \item 端到端重放应用历史并在密码学上链回到 owner。
\end{itemize}

这是 agent-native web。传统 web 应用对人暴露 HTML、对程序暴露文档化
endpoint；邻居程序需要文档、OpenAPI 规约、SDK 与信任。本视图把这四样
合并。

% ──────────────────────────────────────────────────────────────────────────
\section{架构}
\label{sec:architecture}

\subsection{模块拓扑（EasyNet-Cli）}

\begin{lstlisting}
src/runtime/state/                    [与 pages 共享，P-A1]

src/runtime/blob_store/               [与 pages 共享，P-A1]
  tree.rs                             ← 新：TreeBlob put/get、
                                       per-path lookup

src/runtime/agents/webapp_ability.rs  ← 新（P-B2）
                                      register: create / build /
                                      update_config / serve / stop /
                                      delete / get / list

src/runtime/build_executor/           ← 新（P-B1）
  mod.rs                              长任务规范 transition 驱动
  source_snapshot.rs                  项目目录的 Merkle 哈希
  spawn.rs                            子进程 + 日志流
  artifact_capture.rs                 dist 树 -> TreeBlob

src/runtime/views/webapp/             ← 新（P-B5）
  agent_view.rs                       §6 的 JSON 投影
  health.rs
  receipts.rs

src/runtime/hub/                      [扩展 P-A4 的 pages Hub]
  router.rs                           ← 扩展以读取 routes_manifest
  static_handler.rs                   ← 新（P-B3）：TreeBlob 拉取
  api_handler.rs                      ← 新（P-B3）：/api/<verb> 翻译
  ws_handler.rs                       ← 新（P-B4）：/ws/* 升级与
                                       stream 绑定
\end{lstlisting}

\subsection{守护进程接线}

启动时守护进程在 pages 接线之上加：

\begin{enumerate}[leftmargin=2.2em,itemsep=0.15em]
    \item 在 \texttt{\$EASYNET\_HOME/build-cache/} 打开构建执行器工作目录。
    \item 给每个 agent 注册 \texttt{webapp\_ability}。
    \item 对每个规范状态为 \texttt{Building} 的 web\_app（即守护进程
    重启打断了一次构建）：标 Failed 并把 reason 设为 "守护进程在构建
    途中重启"。（按 TR-INV-1 纪律：仅终结规范回执推进状态；中断的
    构建从未触达终结，因此规范状态回滚到 Building 之前的状态——
    \texttt{Created} 或 \texttt{Failed}。重置为 Failed 是\emph{新}
    transition，不是回滚。）
\end{enumerate}

\subsection{Hub 接线}

Hub 二进制 \texttt{easynet-hub} 在 P-B3/P-B4 中加上
\texttt{static\_handler} / \texttt{api\_handler} /
\texttt{ws\_handler}，与 P-A4 的 page handler 并列。其顶层分派为：

\begin{lstlisting}
fn handle(req: HttpRequest) -> HttpResponse {
    let envelope = build_envelope(&req);                  // §5.7 PA
    let (realm, owner, slug, tail) = parse_path(&req.uri)?;
    let kind = state_store.classify(realm, owner, slug)?; // page or web_app

    match kind {
        Kind::Page    => pages_dispatch(envelope, owner, slug, tail),
        Kind::WebApp  => webapp_dispatch(envelope, owner, slug, tail),
        Kind::Unknown => http_404(),
    }
}
\end{lstlisting}

Hub 即使同时服务 pages 与 web\_app 也是单一二进制。path 方案仅按
state 存储中 (\texttt{owner, slug}) 注册的内容来区分两者——URI 形态
完全相同，因为两者都是同一 agent 命名空间下的 state object。

% ──────────────────────────────────────────────────────────────────────────
\section{实现计划}
\label{sec:impl-plan}

五个 PR 切片，可顺序化。切片 P-B1 基于 P-A1（pages 附录的 blob 存储与
state 存储骨架）；P-B3/B4 基于 P-A4（基本 Hub 二进制）。

\subsection{P-B1：构建执行器 + TreeBlob}

\textbf{新增文件。} \texttt{src/runtime/build\_executor/}（4 个模块）、
\texttt{src/runtime/blob\_store/tree.rs}。

\textbf{公开 API。}
\begin{lstlisting}
pub trait BuildExecutor: Send + Sync {
    fn execute(
        &self,
        project_dir: &Path,
        config: &BuildConfig,
        progress: ProgressSink,    // 通往 OperationalReceipt 的 SSE 通道
    ) -> Result<BuildOutcome>;
}
pub enum BuildOutcome {
    Success { artifact_hash: [u8; 32], tree_blob: TreeBlob, log_blob: [u8;32] },
    Failure { reason_hash:   [u8; 32], log_blob:   [u8; 32], exit_status: i32 },
}

pub fn source_snapshot_hash(project_dir: &Path, ignore: &[Pattern])
    -> Result<[u8; 32]>;

pub fn capture_dist(dist_dir: &Path, blob: &dyn BlobStore)
    -> Result<(TreeBlob, [u8; 32])>;
\end{lstlisting}

\textbf{测试。}
\begin{itemize}[leftmargin=2.2em,itemsep=0.1em]
    \item 快照确定性：同一棵树两次快照得相同哈希，与文件遍历顺序无关。
    \item 构建烟测：一个微型 shell 脚本 "构建器"
    （\texttt{mkdir dist \&\& echo hi > dist/index.html}）端到端驱动
    执行器；断言 TreeBlob 往返且产物哈希可确定性重算。
    \item 守护进程在构建中崩溃：启动一个 sleep 的构建；SIGKILL
    守护进程；启动后断言 web\_app 的规范状态被重置为 Failed 并附文档化
    reason。
\end{itemize}

\subsection{P-B2：webapp.* 能力}

\textbf{新增文件。} \texttt{src/runtime/agents/webapp\_ability.rs}。

\textbf{行为。} \S4.5 中每个 transition 都被实现，包括长任务构建
（按 TR-INV-1 流式发出进度 OperationalReceipt 并每次尝试发出且仅发出
一条终结 CanonicalReceipt）。

\textbf{测试。}
\begin{itemize}[leftmargin=2.2em,itemsep=0.1em]
    \item 往返 create/update\_config/build/serve/stop/delete。
    \item 拒绝从 \texttt{Created} 出发的 \texttt{webapp.serve}（规范前置：
    Built）。
    \item 当 \texttt{<verb>} 不在 \texttt{api\_surface} 时拒绝
    \texttt{api/<verb>} 翻译。
    \item 乐观并发：两个针对同一 web\_app 的 \texttt{webapp.update\_config}
    并发尝试；一个赢，一个被 ETag 不一致拒绝。
    \item 构建进度回执可丢失：随机丢一些；重放仍达相同终结规范状态。
\end{itemize}

\subsection{P-B3：Hub 静态 + API 路由}

\textbf{新增文件。} \texttt{src/runtime/hub/static\_handler.rs}、
\texttt{src/runtime/hub/api\_handler.rs}；扩展
\texttt{src/runtime/hub/router.rs} 以读取 \texttt{routes\_manifest}。

\textbf{行为。}
\begin{enumerate}[leftmargin=2.2em,itemsep=0.1em]
    \item 静态处理器：TreeBlob path lookup、blob 拉取、扩展名 MIME
    检测（允许列表）、由 \texttt{csp\_hash} 构 CSP，对 \texttt{/assets/*}
    用 immutable 缓存。
    \item API 处理器：route regex 捕获 verb、按 \texttt{api\_surface}
    校验、把 JSON body 或 query string 解析为 args、调用
    \texttt{<owner>.<verb>}、映射结果。
\end{enumerate}

\textbf{测试。}
\begin{itemize}[leftmargin=2.2em,itemsep=0.1em]
    \item 集成：在临时目录里启动守护进程 + Hub，构建一个 fixture
    web\_app，其 \texttt{api}.abilities 含已注册的 \texttt{shop.echo}
    （回填 args）；\texttt{POST /api/echo \{"x":1\}} 得 \texttt{200 \{"x":1\}}。
    \item 未授权 verb：\texttt{POST /api/secret} 即使 \texttt{<owner>.secret}
    已注册但不在 \texttt{api\_surface} 中也返 404。
    \item 静态 immutable：对 \texttt{/assets/} 路径带 \texttt{If-None-Match}
    两次 GET；第二次响应为 304。
\end{itemize}

\subsection{P-B4：WebSocket 代理 + SPA fallback}

\textbf{新增文件。} \texttt{src/runtime/hub/ws\_handler.rs}。

\textbf{行为。}
\begin{enumerate}[leftmargin=2.2em,itemsep=0.1em]
    \item \texttt{/ws/*} 上 WS 升级；首条文本消息必须是 JSON
    \texttt{\{ability,args\}}；绑到 \texttt{ability stream}；转发后续
    文本帧。
    \item SPA fallback：若 \texttt{static} 或其他规则均不命中，
    用 \texttt{no-cache} 服务 \texttt{dist/index.html}。
\end{enumerate}

\textbf{测试。} WS 流：fixture \texttt{shop.tick} 每 100ms 流式发出
\texttt{\{tick:N\}}；客户端通过 \texttt{/ws/tick} 订阅；断言 600ms 内
收 5 帧；客户端关 WS 后 500ms 内守护进程的流订阅关闭。

\subsection{P-B5：agent\_view + meta 端点}

\textbf{新增文件。} \texttt{src/runtime/views/webapp/}（3 模块）；
对 \texttt{/\_\_easynet/*} 前缀做轻微路由扩展。

\textbf{行为。}
\begin{enumerate}[leftmargin=2.2em,itemsep=0.1em]
    \item \texttt{/\_\_easynet/agent} 返回 \S6 的 JSON。
    \item \texttt{/\_\_easynet/health} 仅返规范与 runtime 状态。
    \item \texttt{/\_\_easynet/receipts} 按 version 顺序返回规范回执，
    可由 query param 约束。
    \item \texttt{/\_\_easynet/schema/<verb>} 返回该 ability 的输入/
    输出 schema。
\end{enumerate}

\textbf{测试。}
\begin{itemize}[leftmargin=2.2em,itemsep=0.1em]
    \item 重放：客户端拉 \texttt{/\_\_easynet/receipts}，按每条回执重构
    规范状态，并计算与 agent\_view 的 \texttt{etag} 匹配的
    \texttt{canonical\_state\_hash}。
    \item Schema：agent\_view 的 \texttt{api.endpoints[*].schema\_url}
    可解析到真实 JSON schema，且该 schema 不拒绝 API 处理器接受的 args。
\end{itemize}

% ──────────────────────────────────────────────────────────────────────────
\section{端到端 worked example}
\label{sec:e2e}

\subsection{准备}

\begin{lstlisting}
$ easynet agent add alice --type claude-code
$ easynet-daemon &
$ easynet-hub --bind 127.0.0.1:8787 &
\end{lstlisting}

\subsection{编辑项目}

Claude Code 以 alice 身份在 alice 的 workspace 中：

\begin{lstlisting}
$ cd $EASYNET_HOME/workspaces/alice
$ npm create vite@latest web_apps/shop -- --template react-ts
$ cd web_apps/shop
$ # ... 编辑 src/App.tsx，写产品列表与"结账"按钮
$ # ... 写 §3.2 中的 easynet.app.toml
\end{lstlisting}

它通过 EasyNet 能力脚手架（不在本附录范围内；见"ability authoring"
附录）注册 API ability：

\begin{lstlisting}
$ easynet ability new shop.list_items --description "列出商品" --lang sh
$ easynet ability new shop.checkout   --description "下单"     --lang sh
$ easynet ability deploy abilities/shop.list_items --node alice
$ easynet ability deploy abilities/shop.checkout   --node alice
\end{lstlisting}

\subsection{创建、构建、上线}

\begin{lstlisting}
$ easynet ability invoke alice.webapp.create --args '{
    "slug":"shop",
    "config_dir":"web_apps/shop"
}'
{ "state":"Created", "version":1, "canonical_state_hash":"0xa3..." }

$ easynet ability invoke alice.webapp.build --args '{ "slug":"shop" }' \
   --stream
[build.started   ] command="npm install && npm run build"
[build.progress  ] log_hash=0x12... 8% installed
[build.progress  ] log_hash=0x13... 47% bundled
[build.progress  ] log_hash=0x14... 92% optimised
[build.completed ] artifact_hash=0x4d28ee17... duration=38.2s
{ "state":"Built", "version":2,
  "build_artifact_hash":"0x4d28ee17..." }

$ easynet ability invoke alice.webapp.serve --args '{ "slug":"shop" }'
{ "runtime_state":"Running" }
\end{lstlisting}

\subsection{浏览器拉取}

\begin{lstlisting}
$ curl -i http://127.0.0.1:8787/default/alice/shop/

HTTP/1.1 200 OK
Content-Type: text/html
ETag: "TSjuFw0pAvjpAQ=="
Cache-Control: max-age=60, must-revalidate
Content-Security-Policy: default-src 'self'; script-src 'self'; ...
Content-Length: 1247

<!doctype html><html><head>...<script type="module" src="/assets/index-DfX9k.js">...
\end{lstlisting}

随后浏览器拉 \texttt{/assets/index-DfX9k.js} ——
\texttt{public, max-age=31536000, immutable}，从 blob 存储服务。
SPA 渲染、列出商品，用户点 Buy。

\subsection{浏览器侧 API 调用}

\begin{lstlisting}
POST /default/alice/shop/api/checkout HTTP/1.1
Content-Type: application/json
Cookie: easynet_session=sess-abc

{ "item_id":"mug-12oz", "qty":1 }

────────────────────────────────────────────────────
HTTP/1.1 200 OK
Content-Type: application/json

{ "order_id":"ord_77a3f", "eta":"2026-04-29" }
\end{lstlisting}

Hub：
\begin{enumerate}[leftmargin=2em,itemsep=0.1em]
    \item 铸造 \texttt{caller=easynet:///principal/human-anon/9c4e.../sess-abc}。
    \item 校验 \texttt{checkout} 在 \texttt{api\_surface} 中（是）。
    \item 用解析后 body 作 args 调用 \texttt{alice.shop.checkout}。
    \item 把回执 result body 返给浏览器。
\end{enumerate}

\subsection{邻居 agent 发现该应用}

\begin{lstlisting}
$ curl http://127.0.0.1:8787/default/alice/shop/__easynet/agent | jq .api
{
  "endpoints": [
    { "verb":"shop.list_items","http":"GET  /api/list_items",
      "schema_url":"/__easynet/schema/shop.list_items" },
    { "verb":"shop.checkout","http":"POST /api/checkout",
      "schema_url":"/__easynet/schema/shop.checkout" }
  ]
}
\end{lstlisting}

邻居 agent——比如运行在另一个节点上的 bob——拉取该视图，读
\texttt{shop.checkout} 的 schema，然后通过 EasyNet 总线（跨节点）
或公开 \texttt{/api/checkout}（同一 Hub）触发该 ability。两条路径
触达同一 ability handler，产生同种 CanonicalReceipt。

\subsection{审计重放}

\texttt{shop} web\_app 的历史显示三条 CanonicalReceipt（create、build、
build）：每条都带 \texttt{post\_state\_hash} 与 \texttt{owner\_signature}。
浏览器驱动的 checkout 触发了一条针对完全不同 state\_type 的
CanonicalReceipt（一个由 alice 或 alice 的订单服务拥有的 \texttt{order}
state object）；该回执也链回 alice 的签名密钥。同时阅读两条链的邻居
可逐字节重构：

\begin{itemize}[leftmargin=2em,itemsep=0.15em]
    \item 当 checkout 发生时，page 处于哪个构建产物；
    \item 那个构建声明了哪些 API endpoint；
    \item checkout 请求触达 \texttt{shop.checkout} 第 K 版并产出订单
    \texttt{ord\_77a3f}；
    \item 上述全部由 alice 签名。
\end{itemize}

\noindent 这是传统 web 决然讲不出的审计故事。今天给 shop 服务的 CDN、
明天的 A/B 重写、第三天的带外订单处理，把真相切成三块；web\_app 把它
合并为一条签名回执链。

% ──────────────────────────────────────────────────────────────────────────
\section{暂缓至 RFC-006-B v2 的开放问题}
\label{sec:open}

\begin{enumerate}[leftmargin=2.2em,itemsep=0.18em]
    \item \textbf{需要持续进程的服务端渲染。} Phase 1 支持
    \texttt{serve.mode = static-only}。需要 Node.js 进程响应每个请求
    的框架（Next.js SSR、Remix loader）需要 \texttt{serve.mode = process}：
    执行器绑定端口、运行进程、Hub 反向代理。runtime 叠加状态那时
    会承载 \texttt{bound\_port}、\texttt{process\_id}、
    \texttt{health\_endpoint}；进程崩溃把 runtime 转入 \texttt{Crashed}
    而不触动 canonical（TR-INV-5 已要求）。

    \item \textbf{跨节点 Hub。} A 节点拥有的 web\_app 由 B 节点上的
    Hub 服务。Hub 升级为联邦客户端；API 处理器经联邦而非本地调用
    \texttt{<owner>.<verb>}。

    \item \textbf{已签资源 bundle 用于 SRI。} TreeBlob 已经按文件提供
    哈希。把这些哈希嵌为 \texttt{index.html} 的 Subresource Integrity
    属性是一个直白的 post-build 阶段；Phase 1 把 SRI 管线留给框架
    自身的插件。

    \item \textbf{构建沙箱。} 见 \S\ref{sec:build-pipeline} 的待决策。

    \item \textbf{多版本路由。} 用户希望 \texttt{/v7/...} 服务构建版本 7，
    而 \texttt{/...} 服务最新版。今天每个 slug 在规范状态里只能有
    单版本。v2 可能引入 \texttt{webapp.tag} transition 把版本与可
    路由别名关联。
\end{enumerate}

\subsection{覆盖检查}

\begin{center}
\renewcommand{\arraystretch}{1.2}
\begin{tabularx}{\linewidth}{l X l}
\toprule
\textbf{主文规则} & \textbf{由谁兑现} & \textbf{位置} \\
\midrule
SO-INV-1（key 不可变）                   & state\_key 解析器拒绝 key 更新                          & \S\ref{sec:state-machine} \\
CSH-INV-1（hash 取自 projection）       & schema 驱动的 $\pi$、JCS、SHA-256                          & \S\ref{sec:state-machine} \\
CSH-INV-2（内容外置为 blob）             & TreeBlob + 单文件 blob；规范状态只持哈希                 & \S\ref{sec:build-pipeline} \\
RC-INV-1（receipt 绑定）                 & 每条 webapp 回执均带 (state\_key, transition\_id, attempt\_id) & \S\ref{sec:state-machine} \\
RC-INV-2（乐观并发）                    & 在 update\_config 上做 ETag 检查                          & \S\ref{sec:impl-plan} P-B2 \\
TA-INV-1（必须带 transition\_id）        & 每个 \texttt{webapp.<verb>@v1} 字符串均被校验             & \S\ref{sec:state-machine} \\
VP-INV-1（ability\_surface 来自 state）  & 由 state\_value 派生 \texttt{ability\_surface}            & \S\ref{sec:agent-view} \\
TR-INV-1（长任务终结）                  & 构建只发恰一条终结 CanonicalReceipt                       & \S\ref{sec:build-pipeline} \\
TR-INV-2（progress 不进规范）           & \texttt{build.progress} 是 OperationalReceipt              & \S\ref{sec:build-pipeline} \\
TR-INV-5（runtime 叠加独立）            & runtime 崩溃不触动 canonical                              & \S\ref{sec:state-machine} \\
TR-INV-7（multi-view 一等公民）          & human（HTML+API+WS）与 agent\_view 都已实现              & \S\ref{sec:hub-routing}, \S\ref{sec:agent-view} \\
TR-INV-8（cross-state-space 矩阵）       & canonical $\times$ runtime 准入被强制                    & \S\ref{sec:state-machine} \\
TR-INV-11（URA-receipt 可追踪）          & \texttt{/\_\_easynet/receipts} + state\_key 索引          & \S\ref{sec:e2e} \\
TR-INV-12（principal/human-anon）        & Hub 在每个请求上铸造该 URI                                & \S\ref{sec:hub-routing} \\
TR-INV-13（caller\_context）             & 由请求头填充；记录到 operational receipt                  & \S\ref{sec:hub-routing}, \S\ref{sec:e2e} \\
WB-INV-1（api\_surface 强制）            & Hub 拒绝未列出 verb 并返 404                              & \S\ref{sec:project} \\
\bottomrule
\end{tabularx}
\end{center}

\paragraph{结语。} pages 是简单的一半：单一文档作为规范状态，无构建、
无 runtime 叠加、无 API。web\_app 是难的一半：长任务规范 transition、
runtime 叠加、声明式 API 面、完整 Hub 路由矩阵、WebSocket。落地
P-B1 至 P-B5 完成有状态 EasyNet web 层 Phase 1，并以工业级形式
落地 agent-native web。

\vfill
\hrule
\vspace{0.4em}
\noindent\small
\textit{历史类比：早期 web 是静态文档（\texttt{easynet.page}）；现代
web 是应用（\texttt{easynet.web\_app}）。EasyNet 规约有意走相同的弧。
先 pages、后 web\_app，每一层都为下一层将要倚仗的协议原语作背书。}

\end{document}
